\documentclass[11pt]{article}
\usepackage{amsmath,amssymb,amsthm}
\usepackage{algorithm,algorithmic}
\usepackage{booktabs}
\usepackage{enumitem}
\usepackage{hyperref}
\usepackage{geometry}
\geometry{margin=1in}
\newtheorem{theorem}{Theorem}
\newtheorem{lemma}{Lemma}
\newtheorem{assumption}{Assumption}
\newcommand{\sign}{\operatorname{sign}}
\newcommand{\EE}{\mathbb{E}}
\newcommand{\PP}{\mathbb{P}}
\newcommand{\RR}{\mathbb{R}}
\newcommand{\norm}[1]{\left\lVert #1\right\rVert}
\begin{document}

\section*{SignSGD (Sign Stochastic Gradient Descent)}

\section*{Mathematical Formulation}
Let $f:\RR^{d}\rightarrow\RR$ be a differentiable (generally non‑convex) objective.  
At iteration $t$ we draw a stochastic gradient $g_t\in\RR^{d}$ (e.g. on a minibatch) and update the parameters $x_t\in\RR^{d}$ by
\[
\boxed{%
x_{t+1}=x_{t}-\alpha\,\sign(g_t),\qquad 
\sign(g_t)_i=
\begin{cases}
+1 & \text{if }[g_t]_i>0,\\[2pt]
-1 & \text{if }[g_t]_i<0,\\[2pt]
0  & \text{if }[g_t]_i=0,
\end{cases}
}
\tag{1}
\]
where $\alpha>0$ is a (scalar) stepsize common to all coordinates.  
The algorithm uses **only the sign** of each stochastic gradient component, thus communicating a single bit per dimension.

\section*{Pseudocode}
\begin{algorithm}[H]
\caption{SignSGD}
\begin{algorithmic}[1]
\STATE \textbf{Input:} stepsize $\alpha>0$, max iterations $T$
\STATE \textbf{Initialize:} $x_0\in\RR^{d}$
\FOR{$t=0,1,\dots,T-1$}
    \STATE Sample a minibatch (or random seed) and compute stochastic gradient $g_t$.
    \STATE $d_t\gets\sign(g_t)$                     \COMMENT{element‑wise sign}
    \STATE $x_{t+1}\gets x_t-\alpha\,d_t$          \COMMENT{update \eqref{1}}
\ENDFOR
\STATE \textbf{Return:} $x_T$
\end{algorithmic}
\end{algorithm}

\section*{Dry Run Example}
Consider the one‑dimensional quadratic
\[
f(w)=(w-3)^2,\qquad \nabla f(w)=2(w-3).
\]
We use a deterministic gradient (so $g_t=\nabla f(w_t)$), stepsize $\alpha=0.1$, and start at $w_0=0$.

\begin{center}
\begin{tabular}{c|c|c|c}
Iteration $t$ & $w_t$ & $g_t= \nabla f(w_t)$ & $w_{t+1}=w_t-\alpha\,\sign(g_t)$\\\midrule
0 & $0$   & $2(0-3)=-6$ & $0-0.1\cdot(-1)=0.1$\\
1 & $0.1$ & $2(0.1-3)=-5.8$ & $0.1-0.1\cdot(-1)=0.2$\\
2 & $0.2$ & $2(0.2-3)=-5.6$ & $0.2-0.1\cdot(-1)=0.3$\\
\end{tabular}
\end{center}
The objective values are $f(0)=9$, $f(0.1)=8.41$, $f(0.2)=7.84$, $f(0.3)=7.29$, showing a steady descent despite using only the sign information.

\newpage
\section*{Assumptions}
\begin{assumption}[Smoothness]\label{ass:smooth}
$f$ has $L$‑Lipschitz continuous gradient, i.e.
\[
\|\nabla f(x)-\nabla f(y)\|\le L\|x-y\|\qquad\forall x,y\in\RR^{d}.
\]
Equivalently,
\[
f(y)\le f(x)+\langle\nabla f(x),y-x\rangle+\frac{L}{2}\|y-x\|^{2}.
\tag{2}
\]
\end{assumption}
\begin{assumption}[Median‑unbiased stochastic gradient]\label{ass:median}
For each coordinate $i\in\{1,\dots,d\}$ and any $x$, the random variable $[g_t(x)]_i$ satisfies
\[
\operatorname{median}\bigl([g_t(x)]_i\bigr)=\bigl[\nabla f(x)\bigr]_i .
\]
Consequently,
\[
p_i(x):=\PP\!\bigl(\sign([g_t(x)]_i)\neq\sign([\nabla f(x)]_i)\bigr)
\le\frac12-\delta,\qquad \text{for some }\delta\in(0,\tfrac12].
\tag{3}
\]
\end{assumption}
\begin{assumption}[Bounded second moment]\label{ass:var}
There exists $G>0$ such that for all $x$,
\[
\EE\!\bigl[\|g_t(x)\|^{2}\mid x_t=x\bigr]\le G^{2}.
\tag{4}
\]
\end{assumption}
All constants $L,\delta,G$ are assumed known to the analyst but not to the algorithm.

\section*{Main Convergence Theorem}
\begin{theorem}[Expected descent for non‑convex SignSGD]\label{thm:main}
Let $\{x_t\}_{t\ge0}$ be generated by \eqref{1} with a constant stepsize
\[
\alpha\le\frac{\delta}{L\,d}.
\tag{5}
\]
Under Assumptions~\ref{ass:smooth}–\ref{ass:var}, for any $T\ge1$,
\[
\frac{1}{T}\sum_{t=0}^{T-1}\EE\!\bigl[\|\nabla f(x_t)\|_{1}\bigr]
\;\le\;
\frac{f(x_0)-f^{\star}}{\alpha\delta T}
\;+\;
\frac{L\alpha d}{2\delta},
\tag{6}
\]
where $f^{\star}=\inf_{x}f(x)$ and $\|\cdot\|_{1}$ denotes the $\ell_{1}$‑norm.
Consequently, choosing
\[
\alpha = \sqrt{\frac{2\,(f(x_0)-f^{\star})}{L d\,T}}
\tag{7}
\]
(which satisfies (5) for $T$ large enough) yields the rate
\[
\frac{1}{T}\sum_{t=0}^{T-1}\EE\!\bigl[\|\nabla f(x_t)\|_{1}\bigr]
\;=\;O\!\bigl(T^{-1/2}\bigr).
\]
\end{theorem}

\section*{Theoretical Properties}
\begin{itemize}[leftmargin=*]
\item \textbf{Stability.} The bound (6) shows that the algorithm is stable provided $\alpha$ respects (5). The term $L\alpha d/2\delta$ originates from the smoothness inequality and vanishes as $\alpha\to0$.
\item \textbf{Hyper‑parameter sensitivity.} Only the scalar stepsize $\alpha$ must be tuned. Because the bound depends on $d$, a smaller $\alpha$ is advisable in high dimensions, matching intuition that a single‑bit update is coarse.
\item \textbf{Comparison to full‑gradient SGD.} For SGD with unbiased gradients one obtains an $O(1/\sqrt{T})$ bound on $\EE\|\nabla f\|^{2}$; SignSGD attains the same $O(T^{-1/2})$ rate on the $\ell_{1}$‑norm, which is comparable up to a factor $d$ (the conversion $\|\cdot\|_{2}\le\|\cdot\|_{1}\le\sqrt{d}\|\cdot\|_{2}$).
\item \textbf{Median‑based variance bound.} Condition (3) is weaker than the usual unbiasedness assumption; it only requires that the sign is correct with probability $>1/2$, a property that holds for many heavy‑tailed noise models.
\end{itemize}

\section*{Discussion}
The theorem demonstrates that even though SignSGD discards magnitude information, the algorithm still converges to a stationary point in expectation for smooth non‑convex objectives. The use of the median‑unbiased assumption replaces the classic variance bound and yields the clean probability $p_i\le 1/2-\delta$ that appears directly in the descent inequality. The proof below follows a first‑principles expansion of the smoothness inequality, careful conditioning on the stochastic gradient, and an explicit summation over coordinates.

\newpage
\section*{Preliminaries (Definitions and Lemmas)}
\begin{lemma}[Coordinate‑wise sign error probability]\label{lem:signerror}
For any $x\in\RR^{d}$ and any coordinate $i$, let $p_i(x)$ be defined as in \eqref{3}. Then
\[
\EE\!\bigl[\sign([g_t(x)]_i)\mid x_t=x\bigr]
= (1-2p_i(x))\,\sign([\nabla f(x)]_i).
\tag{8}
\]
\end{lemma}
\begin{proof}
Condition on the event that the sign matches or mismatches:
\[
\EE[\sign(g_{t,i})]=
\PP(\text{match})\cdot\sign(\nabla f_i)
+
\PP(\text{mismatch})\cdot(-\sign(\nabla f_i))
=
(1-p_i-p_i)\sign(\nabla f_i)
=(1-2p_i)\sign(\nabla f_i).
\]
\end{proof}

\begin{lemma}[Bound on squared sign vector]\label{lem:squared}
For any stochastic gradient $g_t$,
\[
\|\sign(g_t)\|^{2}=d .
\tag{9}
\]
\end{lemma}
\begin{proof}
Each coordinate of $\sign(g_t)$ belongs to $\{-1,0,+1\}$, and by definition at least one of $\pm1$ occurs unless the coordinate is exactly zero (probability zero for continuous distributions). Hence the square of each entry is $1$, and summing over $d$ coordinates gives $d$.
\end{proof}

\newpage
\section*{Main Proof (Step‑by‑Step Derivation)}
We prove Theorem~\ref{thm:main}. Throughout the proof we condition on the sigma‑algebra $\mathcal{F}_t$ generated by $\{x_0,\dots,x_t\}$.

\noindent\textbf{Step 1 – Smoothness expansion.}  
From the $L$‑smoothness inequality \eqref{2} with $y=x_{t+1}=x_t-\alpha\sign(g_t)$ we obtain
\[
f(x_{t+1})\le f(x_t)-\alpha\langle\nabla f(x_t),\sign(g_t)\rangle
+\frac{L\alpha^{2}}{2}\|\sign(g_t)\|^{2}.
\tag{10}
\]
\noindent\textbf{Step 2 – Take conditional expectation.}  
Apply $\EE[\cdot\mid\mathcal{F}_t]$ to \eqref{10} and use Lemma~\ref{lem:squared}:
\[
\EE\!\bigl[f(x_{t+1})\mid\mathcal{F}_t\bigr]
\le f(x_t)-\alpha\EE\!\bigl[\langle\nabla f(x_t),\sign(g_t)\rangle\mid\mathcal{F}_t\bigr]
+\frac{L\alpha^{2}d}{2}.
\tag{11}
\]
\noindent\textbf{Step 3 – Expand the inner product coordinate‑wise.}  
\[
\langle\nabla f(x_t),\sign(g_t)\rangle
=\sum_{i=1}^{d}[\nabla f(x_t)]_i\,\sign([g_t]_i).
\tag{12}
\]
\noindent\textbf{Step 4 – Compute the conditional expectation of each term.}  
Using Lemma~\ref{lem:signerror} and the definition $p_i:=p_i(x_t)$,
\[
\EE\!\bigl[\sign([g_t]_i)\mid\mathcal{F}_t\bigr]
=(1-2p_i)\,\sign([\nabla f(x_t)]_i).
\tag{13}
\]
Hence,
\[
\EE\!\bigl[[\nabla f(x_t)]_i\sign([g_t]_i)\mid\mathcal{F}_t\bigr]
=|\,[\nabla f(x_t)]_i\,|\,(1-2p_i).
\tag{14}
\]
\noindent\textbf{Step 5 – Lower bound the expected descent term.}  
Since $p_i\le\frac12-\delta$ by Assumption~\ref{ass:median},
\[
1-2p_i\ge 2\delta.
\tag{15}
\]
Insert (14) into (11) and use (15):
\[
\EE\!\bigl[f(x_{t+1})\mid\mathcal{F}_t\bigr]
\le f(x_t)-2\alpha\delta\sum_{i=1}^{d}|\,[\nabla f(x_t)]_i\,|
+\frac{L\alpha^{2}d}{2}.
\tag{16}
\]
\noindent\textbf{Step 6 – Rewrite with $\ell_{1}$‑norm.}  
Define $\|\nabla f(x_t)\|_{1}=\sum_{i=1}^{d}|\,[\nabla f(x_t)]_i\,|$. Then (16) becomes
\[
\EE\!\bigl[f(x_{t+1})\mid\mathcal{F}_t\bigr]
\le f(x_t)-2\alpha\delta\|\nabla f(x_t)\|_{1}
+\frac{L\alpha^{2}d}{2}.
\tag{17}
\]
\noindent\textbf{Step 7 – Unconditional expectation.}  
Taking total expectation and rearranging,
\[
2\alpha\delta\;\EE\!\bigl[\|\nabla f(x_t)\|_{1}\bigr]
\le \EE\!\bigl[f(x_t)-f(x_{t+1})\bigr]
+\frac{L\alpha^{2}d}{2}.
\tag{18}
\]
\noindent\textbf{Step 8 – Sum over $t=0,\dots,T-1$.}  
Summing (18) telescopically yields
\[
2\alpha\delta\sum_{t=0}^{T-1}\EE\!\bigl[\|\nabla f(x_t)\|_{1}\bigr]
\le f(x_0)-\EE\!\bigl[f(x_T)\bigr]
+\frac{L\alpha^{2}dT}{2}.
\tag{19}
\]
Since $f(x_T)\ge f^{\star}$, we replace $\EE[f(x_T)]$ by $f^{\star}$ to obtain an upper bound:
\[
2\alpha\delta\sum_{t=0}^{T-1}\EE\!\bigl[\|\nabla f(x_t)\|_{1}\bigr]
\le f(x_0)-f^{\star}
+\frac{L\alpha^{2}dT}{2}.
\tag{20}
\]
\noindent\textbf{Step 9 – Divide by $2\alpha\delta T$ to obtain the average.}
\[
\frac{1}{T}\sum_{t=0}^{T-1}\EE\!\bigl[\|\nabla f(x_t)\|_{1}\bigr]
\le
\frac{f(x_0)-f^{\star}}{2\alpha\delta T}
+\frac{L\alpha d}{4\delta}.
\tag{21}
\]
Multiplying numerator and denominator of the first term by $2$ gives exactly the statement (6) of Theorem~\ref{thm:main}.

\noindent\textbf{Step 10 – Choice of stepsize.}  
Set $\alpha$ as in \eqref{7}. Substituting into (21) yields
\[
\frac{1}{T}\sum_{t=0}^{T-1}\EE\!\bigl[\|\nabla f(x_t)\|_{1}\bigr]
\le
\frac{2\sqrt{(f(x_0)-f^{\star})L d}}{2\delta\sqrt{T}}
+\frac{L d}{4\delta}\sqrt{\frac{2(f(x_0)-f^{\star})}{L d T}}
=O\!\bigl(T^{-1/2}\bigr).
\tag{22}
\]
Thus the claimed $O(T^{-1/2})$ rate follows.

\noindent\textbf{Step 11 – Verification of stepsize condition (5).}  
With $\alpha$ from \eqref{7},
\[
\alpha = \sqrt{\frac{2\,(f(x_0)-f^{\star})}{L d\,T}}
\le \frac{\delta}{L d}
\quad\Longleftrightarrow\quad
T\ge \frac{2\,L d\,(f(x_0)-f^{\star})}{\delta^{2}}.
\]
Hence for any $T$ satisfying the right‑hand side, the stepsize respects (5) and all previous inequalities are valid. \hfill $\square$

\section*{Rate Derivation (With Constants)}
From (21) we have the explicit bound
\[
\frac{1}{T}\sum_{t=0}^{T-1}\EE\!\bigl[\|\nabla f(x_t)\|_{1}\bigr]
\le
\underbrace{\frac{f(x_0)-f^{\star}}{2\delta}}_{C_{1}}
\cdot\frac{1}{\alpha T}
\;+\;
\underbrace{\frac{L d}{4\delta}}_{C_{2}}\alpha .
\tag{23}
\]
Minimising the right‑hand side over $\alpha>0$ yields $\alpha^{\star}= \sqrt{\frac{2(f(x_0)-f^{\star})}{L d T}}$ and the optimal bound
\[
\frac{1}{T}\sum_{t=0}^{T-1}\EE\!\bigl[\|\nabla f(x_t)\|_{1}\bigr]
\le
\frac{\sqrt{2 L d (f(x_0)-f^{\star})}}{2\delta}\,\frac{1}{\sqrt{T}}.
\tag{24}
\]
Thus the constant in the $O(T^{-1/2})$ rate is $\frac{\sqrt{2 L d (f(x_0)-f^{\star})}}{2\delta}$.

\section*{Special Cases}
\begin{enumerate}[leftmargin=*]
\item \textbf{Exact gradients.} If $g_t=\nabla f(x_t)$ deterministically, then $p_i=0$ and $\delta=1/2$. Inequality (17) becomes
\[
\EE[f(x_{t+1})]\le f(x_t)-\alpha\|\nabla f(x_t)\|_{1}+\frac{L\alpha^{2}d}{2},
\]
recovering the classic deterministic sign‑descent bound.
\item \textbf{Heavy‑tailed noise.} Even when $\EE[\|g_t\|^{2}]$ is infinite, Assumption~\ref{ass:median} may still hold (median is robust). The proof only uses the sign‑error probability $p_i$, so the convergence guarantee remains valid.
\item \textbf{Coordinate‑wise adaptive stepsizes.} If one replaces the scalar $\alpha$ by a vector $\alpha_i$ satisfying $\alpha_i\le \frac{\delta}{L}$ for each $i$, the same derivation yields a bound with $\sum_i \alpha_i$ in place of $\alpha d$.
\end{enumerate}

\end{document}